%
% audio-programs.tex
%
% Z Specification for the Vox "Audio Programs" lifecycle.
% Phase 1 (playlist format) modelled in full, with the shared Program state
% and mode vocabulary that podcast and audiobook formats extend later.
%
% Generalises the earlier four-state playlist model (vox-1rxb) and folds in the failure /
% observability state machine of docs/vox-ig52-music-resilience.md.
%

\documentclass[a4paper,10pt,fleqn]{article}
\usepackage[margin=1in]{geometry}
\usepackage{fuzz}
\usepackage[colorlinks=true,linkcolor=blue,citecolor=blue,urlcolor=blue]{hyperref}

\begin{document}

\title{Vox Audio Programs: A Z Specification}
\author{Formal Model of the Program lifecycle --- Phase 1 (playlist)}
\date{July 2026}
\hypersetup{
  pdftitle={Vox Audio Programs: A Z Specification},
  pdfauthor={Punt Labs},
  pdfsubject={Z Specification},
  pdfcreator={fuzz/probcli}
}
\maketitle

\tableofcontents
\newpage

\section{Introduction}

A \emph{Program} is a named, ordered collection of \emph{Parts}: each Part is
LLM-authored content plus generated audio. The three formats --- \emph{playlist},
\emph{podcast}, \emph{audiobook} --- share one lifecycle: an agent authors a
part, the daemon generates it in the background (which costs credits), the part
is stored, and thereafter it plays, advances, and rotates for free. This
specification pins that lifecycle for Phase~1, in which only the \emph{playlist}
format exists, while fixing the shared state and mode vocabulary so that the
later formats extend the same model rather than replacing it.

The model does two things the informal concept could not. First, it
generalises the earlier four-state playlist machine (vox-1rxb)
into a six-mode Program machine that carries the resilience states
(\texttt{retrying}, \texttt{failed}) from \texttt{vox-ig52}, so that a
generation failure is a \emph{modelled, observable} transition rather than a
silent disable. Second, it adds the Phase-1 consume-only operations ---
\texttt{PlayPart}, \texttt{Loop}, \texttt{StartFromDisk} --- that let the CLI
drive an existing pool with no LLM and zero credits, and it proves, as
preconditions and framing, that those operations generate nothing.

Writing the model surfaced several transitions the concept left underspecified
or contradictory; they are collected in Section~\ref{sec:findings}, which is
the highest-value output of the exercise.

\section{Basic Types}

A \emph{Part} is an opaque saved-audio identifier (in the implementation, the
content-addressed file within a Program's directory). A \emph{reason} is a
human-readable diagnostic string; the model treats it as opaque, but every
inhabitant of \texttt{REASON} is understood to be non-empty (an empty reason is
not a reason).

There is no session or owner type. The daemon's Program state is
machine-universal: any client --- an MCP agent session or the CLI, from any
process --- may issue any command against it. An earlier draft gated control on
the session that turned a Program on; the operator ruled that out
(2026-07-05) as complex, bug-ridden, and unusable across sessions, so the model
carries no notion of ownership at all.

\begin{zed}
  [PART, REASON]
\end{zed}

\section{Free Types}

Three formats share one Program model. Only \texttt{playlist} is realised in
Phase~1; the other two are named so that \texttt{poolSize} and the operations
that branch on it are total from the start.

\begin{zed}
  Format ::= playlist | podcast | audiobook
\end{zed}

The Program mode is a six-state machine. The first four are the playlist modes
of the earlier playlist model; \texttt{retrying} and \texttt{failed} are the
resilience states of \texttt{vox-ig52}. The modes are format-general: for a
finite format (podcast, audiobook) \texttt{playingRotating} denotes ``every
Part generated, playing sequentially'', and the loop-versus-stop distinction at
the end lives in the advance operation, not in a new mode.

\begin{zed}
  Mode ::= off | generatingFirst | playingFilling | playingRotating \\
  \t2 | retrying | failed
\end{zed}

A Part passes through a lifecycle of its own. In Phase~1 a stored Part is either
\texttt{ready} (in the pool, playable) or \texttt{failedPart} (a permanent
generation error, with a retained reason); \texttt{pending} and
\texttt{generating} are in-flight statuses that the atomic-delivery abstraction
below represents by the \texttt{filling} flag rather than by a per-Part stored
field. The free type is declared in full because the spoken formats, whose
generation is long, will promote \texttt{pending}/\texttt{generating} to stored
per-Part state.

\begin{zed}
  PartStatus ::= pending | generating | ready | failedPart
\end{zed}

A Z boolean (avoids the ProB/B keyword clash with \texttt{BOOL}).

\begin{zed}
  ZBOOL ::= ztrue | zfalse
\end{zed}

\section{Global Constants}

\texttt{poolSize} fixes, per format, the number of Parts a full Program holds;
for \texttt{playlist} the design value is 12. The podcast and audiobook values
are provisional placeholders (concept: $\sim$6 chapters, a finite episode run);
Phase~2/3 pin them, and nothing in Phase~1 exercises them. \texttt{maxRetry}
bounds the transient-error backoff before an \emph{empty-pool} Program gives up.

\begin{axdef}
  poolSize : Format \fun \nat_1 \\
  maxRetry : \nat_1
  \where
  poolSize~playlist = 12 \\
  poolSize~podcast = 6 \\
  poolSize~audiobook = 6 \\
  maxRetry = 5
\end{axdef}

\section{State}
\label{sec:state}

\texttt{playing} and \texttt{lastPlayed} are ``optional Parts'' (sets of at most
one Part): empty when nothing is playing / no Part has yet been played.
\texttt{lastError} is an optional reason. \texttt{pool} is the set of
\emph{ready} Parts --- the playable audio on
disk --- and \texttt{failedParts} records the Parts that hit a permanent
generation error together with the retained reason; the two are disjoint (a
Part is ready \emph{or} failed, never both). \texttt{attempts} counts transient
retries in flight.

The predicate carries every invariant. The mode implications
(\texttt{S10}--\texttt{S16} in the prose below) tie each mode to the shape of
the rest of the state, so that ``a full pool never hard-fails'' and ``entering
\texttt{failed} sets an observable reason'' are consequences of the state
schema, not of scattered defensive checks.

\begin{schema}{Program}
  format : Format \\
  pool : \power PART \\
  failedParts : PART \pfun REASON \\
  playing : \power PART \\
  lastPlayed : \power PART \\
  mode : Mode \\
  filling : ZBOOL \\
  attempts : \nat \\
  lastError : \power REASON
  \where
  \# pool \leq (poolSize~format) \\
  \# playing \leq 1 \land \# lastPlayed \leq 1 \\
  \# lastError \leq 1 \\
  playing \subseteq pool \land lastPlayed \subseteq pool \\
  (\dom failedParts) \cap pool = \emptyset \\
  attempts \leq maxRetry \\
  (attempts \geq 1 \iff mode = retrying) \\
  filling = ztrue \implies \\
  \quad~ (mode = generatingFirst \lor mode = playingFilling) \\
  \# lastError = 1 \implies (mode = retrying \lor mode = failed) \\
  mode = failed \implies \# lastError = 1 \\
  mode = off \implies \\
  \quad~ (playing = \emptyset \land filling = zfalse \\
  \quad~ \land lastError = \emptyset \land failedParts = \emptyset) \\
  mode = generatingFirst \implies \\
  \quad~ (pool = \emptyset \land playing = \emptyset \land filling = ztrue) \\
  mode = playingFilling \implies \\
  \quad~ (\# playing = 1 \land \# pool \geq 1 \land \# pool < (poolSize~format)) \\
  mode = playingRotating \implies \\
  \quad~ (\# playing = 1 \land \# pool = (poolSize~format) \land filling = zfalse) \\
  mode = retrying \implies \\
  \quad~ (filling = zfalse \land (pool = \emptyset \iff playing = \emptyset)) \\
  mode = failed \implies \\
  \quad~ (pool = \emptyset \land playing = \emptyset \land filling = zfalse)
\end{schema}

The ready pool and the failed set are the Phase-1 realisation of the Part
status of Section~3: $p \in pool$ means $status(p) = ready$, and
$p \in \dom failedParts$ means $status(p) = failedPart$ with reason
$failedParts(p)$. The invariant $(\dom failedParts) \cap pool = \emptyset$ is
exactly ``a Part is ready or failed, never both''.

\section{Initialisation}

The daemon starts idle with an empty \texttt{playlist} pool --- the only format
realised in Phase~1. \texttt{RestartFromDisk} models a restart in which a pool
of already-saved ready Parts is loaded --- still idle until turned on --- which
is the state from which the CLI plays existing audio with no generation; it is
parameterised by \texttt{fmt?} to show that the initialisation rule is
format-general.

\begin{schema}{Init}
  Program'
  \where
  format' = playlist \\
  pool' = \emptyset \\
  failedParts' = \emptyset \\
  playing' = \emptyset \\
  lastPlayed' = \emptyset \\
  mode' = off \\
  filling' = zfalse \\
  attempts' = 0 \\
  lastError' = \emptyset
\end{schema}

\begin{schema}{RestartFromDisk}
  Program' \\
  fmt? : Format \\
  diskPool? : \power PART
  \where
  \# diskPool? \leq (poolSize~fmt?) \\
  format' = fmt? \\
  pool' = diskPool? \\
  failedParts' = \emptyset \\
  playing' = \emptyset \\
  lastPlayed' = \emptyset \\
  mode' = off \\
  filling' = zfalse \\
  attempts' = 0 \\
  lastError' = \emptyset
\end{schema}

\section{Generation Operations}

These operations run on the MCP path, where an agent authors content and the
daemon spends credits generating it. They spend credits; the consume-only path
(Section~8) does not. Neither path is gated on a session --- any client may
drive either.

\subsection{TurnOn}

From \texttt{off}, a client enters the pool. An empty pool begins generating the
first Part; a partial pool starts playing and resumes the background fill; a
full pool starts playing and rotates.

\begin{schema}{TurnOn}
  \Delta Program
  \where
  mode = off \\
  format' = format \\
  pool' = pool \\
  failedParts' = \emptyset \\
  lastPlayed' = \emptyset \\
  attempts' = 0 \\
  lastError' = \emptyset \\
  (pool = \emptyset \implies \\
    \quad~ (mode' = generatingFirst \land playing' = \emptyset \land filling' = ztrue)) \\
  (\# pool \geq 1 \land \# pool < (poolSize~format) \implies \\
    \quad~ (mode' = playingFilling \land playing' \subseteq pool \\
    \quad~ \land \# playing' = 1 \land filling' = ztrue)) \\
  (\# pool = (poolSize~format) \implies \\
    \quad~ (mode' = playingRotating \land playing' \subseteq pool \\
    \quad~ \land \# playing' = 1 \land filling' = zfalse))
\end{schema}

\subsection{FirstTrackOk}

The generation of the first Part succeeds. It becomes the sole ready Part and
starts playing; the fill continues (unless the format's pool is already full at
one Part).

\begin{schema}{FirstTrackOk}
  \Delta Program \\
  new? : PART
  \where
  mode = generatingFirst \\
  new? \notin \dom failedParts \\
  format' = format \\
  pool' = \{ new? \} \\
  failedParts' = failedParts \\
  playing' = \{ new? \} \\
  lastPlayed' = \emptyset \\
  attempts' = 0 \\
  lastError' = \emptyset \\
  ((poolSize~format) = 1 \implies (mode' = playingRotating \land filling' = zfalse)) \\
  ((poolSize~format) > 1 \implies (mode' = playingFilling \land filling' = ztrue))
\end{schema}

\subsection{FirstTrackBadPrompt}

The first Part hits a \emph{permanent} error (a \texttt{bad\_prompt} / ToS
rejection, or a missing key). There is nothing playing and nothing to fall back
to, so the Program enters \texttt{failed}: the attempted Part is recorded in
\texttt{failedParts} with its reason, and \texttt{lastError} is set. This is an
\emph{observable} transition --- \texttt{mode} flips and a reason appears ---
never a silent disable.

\begin{schema}{FirstTrackBadPrompt}
  \Delta Program \\
  bad? : PART \\
  reason? : REASON
  \where
  mode = generatingFirst \\
  bad? \notin \dom failedParts \\
  format' = format \\
  pool' = \emptyset \\
  failedParts' = failedParts \oplus \{ bad? \mapsto reason? \} \\
  playing' = \emptyset \\
  lastPlayed' = \emptyset \\
  mode' = failed \\
  filling' = zfalse \\
  attempts' = 0 \\
  lastError' = \{ reason? \}
\end{schema}

\subsection{FirstTrackTransient}

The first Part hits a \emph{transient} error (429, quota, 5xx, timeout). The
Program backs off in \texttt{retrying} with an empty pool; the reason is
retained as advisory. No Part playing yet, so \texttt{RetryExhausted} may later
force \texttt{failed}.

\begin{schema}{FirstTrackTransient}
  \Delta Program \\
  reason? : REASON
  \where
  mode = generatingFirst \\
  format' = format \\
  pool' = pool \\
  failedParts' = failedParts \\
  playing' = \emptyset \\
  lastPlayed' = \emptyset \\
  mode' = retrying \\
  filling' = zfalse \\
  attempts' = 1 \\
  lastError' = \{ reason? \}
\end{schema}

\subsection{FillOk}

The background fill delivers one fresh ready Part. Enabled only while a fill is
actually running (\texttt{filling} true) and the pool is not yet full; the Part
must be new (the deterministic collision-free naming guarantees this). Reaching
\texttt{poolSize} stops the fill and switches to rotation.

\begin{schema}{FillOk}
  \Delta Program \\
  new? : PART
  \where
  mode = playingFilling \\
  filling = ztrue \\
  new? \notin pool \\
  new? \notin \dom failedParts \\
  format' = format \\
  pool' = pool \cup \{ new? \} \\
  failedParts' = failedParts \\
  playing' = playing \\
  lastPlayed' = lastPlayed \\
  attempts' = 0 \\
  lastError' = \emptyset \\
  (\# pool' = (poolSize~format) \implies \\
    \quad~ (mode' = playingRotating \land filling' = zfalse)) \\
  (\# pool' < (poolSize~format) \implies \\
    \quad~ (mode' = playingFilling \land filling' = ztrue))
\end{schema}

\subsection{FillBadPart}

A \emph{permanent} error on a background Part while the Program is already
playing. The offending Part is recorded in \texttt{failedParts}, but the pool
is not burned and playback continues: the Program stays healthy
(\texttt{lastError} empty) and keeps filling. The failure is observable
\emph{per Part} --- via \texttt{failedParts} --- not as a program-level error.
This is the second observability surface (finding~\#5).

\begin{schema}{FillBadPart}
  \Delta Program \\
  bad? : PART \\
  reason? : REASON
  \where
  mode = playingFilling \\
  filling = ztrue \\
  bad? \notin pool \\
  bad? \notin \dom failedParts \\
  format' = format \\
  pool' = pool \\
  failedParts' = failedParts \oplus \{ bad? \mapsto reason? \} \\
  playing' = playing \\
  lastPlayed' = lastPlayed \\
  mode' = playingFilling \\
  filling' = ztrue \\
  attempts' = 0 \\
  lastError' = \emptyset
\end{schema}

\subsection{FillTransient}

A \emph{transient} error on a background Part while playing. The Program backs
off in \texttt{retrying} but keeps playing the existing pool (which is
non-empty), so it can never fall through to \texttt{failed} (finding~\#4).

\begin{schema}{FillTransient}
  \Delta Program \\
  reason? : REASON
  \where
  mode = playingFilling \\
  filling = ztrue \\
  format' = format \\
  pool' = pool \\
  failedParts' = failedParts \\
  playing' = playing \\
  lastPlayed' = lastPlayed \\
  mode' = retrying \\
  filling' = zfalse \\
  attempts' = 1 \\
  lastError' = \{ reason? \}
\end{schema}

\subsection{RetryFails}

Another transient error fires during backoff, below the cap. Only the counter
and the advisory reason move; the pool and playback are untouched.

\begin{schema}{RetryFails}
  \Delta Program \\
  reason? : REASON
  \where
  mode = retrying \\
  attempts < maxRetry \\
  format' = format \\
  pool' = pool \\
  failedParts' = failedParts \\
  playing' = playing \\
  lastPlayed' = lastPlayed \\
  mode' = retrying \\
  filling' = zfalse \\
  attempts' = attempts + 1 \\
  lastError' = \{ reason? \}
\end{schema}

\subsection{RetryExhausted}

The backoff cap is reached \emph{and the pool is empty} --- there is nothing to
play and no prospect of a Part --- so the Program gives up into \texttt{failed}
with an observable reason. The empty-pool precondition is what confines
hard-failure to the first-track case: a non-empty pool has no
\texttt{RetryExhausted} enabled (finding~\#4).

\begin{schema}{RetryExhausted}
  \Delta Program \\
  reason? : REASON
  \where
  mode = retrying \\
  attempts = maxRetry \\
  pool = \emptyset \\
  format' = format \\
  pool' = \emptyset \\
  failedParts' = failedParts \\
  playing' = \emptyset \\
  lastPlayed' = \emptyset \\
  mode' = failed \\
  filling' = zfalse \\
  attempts' = 0 \\
  lastError' = \{ reason? \}
\end{schema}

\subsection{RetryCapped}

A further transient error fires during backoff, but the counter is already at
\texttt{maxRetry} and the pool is non-empty. There is something to play, so the
Program neither hard-fails --- \texttt{RetryExhausted} needs an empty pool ---
nor climbs past the cap --- \texttt{RetryFails} is guarded below it. The retry
self-loops in \texttt{retrying}, pinning \texttt{attempts} at the cap and
refreshing only the advisory reason; the pool and playback are untouched. This
is the cap-and-hold that makes ``a non-empty pool tolerates transient errors
indefinitely'' (finding~\#4) a fact of the model rather than a runtime hope:
because \texttt{mode} stays \texttt{retrying}, the invariant
$mode = failed \implies pool = \emptyset$ is never engaged, so a Program with a
non-empty pool cannot fall through to \texttt{failed} at the cap.

\begin{schema}{RetryCapped}
  \Delta Program \\
  reason? : REASON
  \where
  mode = retrying \\
  attempts = maxRetry \\
  pool \neq \emptyset \\
  format' = format \\
  pool' = pool \\
  failedParts' = failedParts \\
  playing' = playing \\
  lastPlayed' = lastPlayed \\
  mode' = retrying \\
  filling' = zfalse \\
  attempts' = attempts \\
  lastError' = \{ reason? \}
\end{schema}

\subsection{Recover}

The backoff clears and generation resumes. From an empty pool the Program
returns to \texttt{generatingFirst} (retry the first Part); from a non-empty
pool it returns to \texttt{playingFilling} with playback intact. In both cases
the advisory reason and the counter are cleared and the fill is re-armed.

\begin{schema}{Recover}
  \Delta Program
  \where
  mode = retrying \\
  format' = format \\
  pool' = pool \\
  failedParts' = failedParts \\
  lastPlayed' = lastPlayed \\
  attempts' = 0 \\
  lastError' = \emptyset \\
  filling' = ztrue \\
  (pool = \emptyset \implies (mode' = generatingFirst \land playing' = \emptyset)) \\
  (pool \neq \emptyset \implies (mode' = playingFilling \land playing' = playing))
\end{schema}

\subsection{VibeStyleChange}

A client retunes to a new \emph{(vibe, style)} key, whose already-saved ready
Parts are \texttt{newPool?}. The selection of mode/playing/filling mirrors
\texttt{TurnOn} applied to the new pool; the old failed Parts, advisory reason,
and avoid-repeat key are discarded. Retuning is allowed only while
the Program is not \texttt{off} --- a forwarded vibe while off is refused, so it
spends no generation credits. Reaching this operation from \texttt{failed}
(re-authoring after a \texttt{bad\_prompt}) is the intended recovery path, and
because there is no ownership it is available to any client.

\begin{schema}{VibeStyleChange}
  \Delta Program \\
  newPool? : \power PART
  \where
  mode \neq off \\
  \# newPool? \leq (poolSize~format) \\
  format' = format \\
  pool' = newPool? \\
  failedParts' = \emptyset \\
  lastPlayed' = \emptyset \\
  attempts' = 0 \\
  lastError' = \emptyset \\
  (newPool? = \emptyset \implies \\
    \quad~ (mode' = generatingFirst \land playing' = \emptyset \land filling' = ztrue)) \\
  (\# newPool? \geq 1 \land \# newPool? < (poolSize~format) \implies \\
    \quad~ (mode' = playingFilling \land playing' \subseteq newPool? \\
    \quad~ \land \# playing' = 1 \land filling' = ztrue)) \\
  (\# newPool? = (poolSize~format) \implies \\
    \quad~ (mode' = playingRotating \land playing' \subseteq newPool? \\
    \quad~ \land \# playing' = 1 \land filling' = zfalse))
\end{schema}

\subsection{TurnOff}

Cancel the fill and stop playback. Every generation-related field is cleared,
so a later forwarded vibe is refused while \texttt{off}. The ready pool remains
on disk for the CLI to play.

\begin{schema}{TurnOff}
  \Delta Program
  \where
  mode \neq off \\
  format' = format \\
  pool' = pool \\
  failedParts' = \emptyset \\
  playing' = \emptyset \\
  lastPlayed' = \emptyset \\
  mode' = off \\
  filling' = zfalse \\
  attempts' = 0 \\
  lastError' = \emptyset
\end{schema}

\section{Consume-only Operations (no generation)}

These operations drive an existing pool. None of them generates a Part or
spends a credit: each leaves \texttt{pool}, \texttt{failedParts}, and
\texttt{filling} exactly as it found them (or, for \texttt{StartFromDisk},
leaves \texttt{filling} false). They only move the playback cursor over ready
Parts, which spends nothing and changes no \emph{(vibe, style)} key. Like every
other operation, they are gated on \texttt{mode} alone, never on a session ---
the CLI can play a pool it did not turn on, which is the common case after a
daemon restart.

\subsection{Rotate --- auto-advance, loop, next, and skip}

When the current Part ends, move to another ready Part, avoiding an immediate
repeat when the pool holds more than one. With a single Part it replays --- the
only intended looping. \texttt{Rotate} is enabled in every playing mode,
\emph{including} \texttt{retrying}: during a transient backoff the existing pool
must keep playing and advancing (finding~\#3). Generation state
(\texttt{mode}, \texttt{filling}, \texttt{attempts}, \texttt{lastError}) is
untouched.

\begin{schema}{Rotate}
  \Delta Program
  \where
  (mode = playingFilling \lor mode = playingRotating \lor mode = retrying) \\
  \# pool \geq 1 \\
  format' = format \\
  pool' = pool \\
  failedParts' = failedParts \\
  mode' = mode \\
  filling' = filling \\
  attempts' = attempts \\
  lastError' = lastError \\
  lastPlayed' = playing \\
  (\# pool \geq 2 \implies \\
    \quad~ (playing' \subseteq pool \land \# playing' = 1 \land playing' \neq playing)) \\
  (\# pool = 1 \implies playing' = playing)
\end{schema}

The automatic track-end advance, the CLI \texttt{loop}/\texttt{next}, and an
explicit \texttt{skip} are all the same transition. With ownership removed,
there is no owner-gated variant: an automatic advance and a user-driven advance
are indistinguishable, and any client may drive either. \texttt{Skip},
\texttt{Next}, \texttt{Loop}, and \texttt{TrackEnds} are therefore all names for
\texttt{Rotate}.

\begin{zed}
  TrackEnds \defs Rotate
\end{zed}

\begin{zed}
  Loop \defs Rotate
\end{zed}

\begin{zed}
  Next \defs Rotate
\end{zed}

\begin{zed}
  Skip \defs Rotate
\end{zed}

\subsection{PlayPart --- replay a specific Part by index}

The CLI selects a specific ready Part (\texttt{vox music playlist playlist:2}).
Unlike \texttt{Rotate}, there is \emph{no} anti-repeat constraint: the user
asked for that Part, so it plays even if it equals the current or just-played
Part. The precondition $target? \in pool$ is exactly ``the addressed index must
resolve to a \texttt{ready} Part''; the CLI resolves the 1-based index to a Part
against a stable numbering of \texttt{pool} (finding~\#7). Nothing is generated.

\begin{schema}{PlayPart}
  \Delta Program \\
  target? : PART
  \where
  (mode = playingFilling \lor mode = playingRotating \lor mode = retrying) \\
  target? \in pool \\
  format' = format \\
  pool' = pool \\
  failedParts' = failedParts \\
  mode' = mode \\
  filling' = filling \\
  attempts' = attempts \\
  lastError' = lastError \\
  lastPlayed' = playing \\
  playing' = \{ target? \}
\end{schema}

\subsection{StartFromDisk --- CLI cold play with no generation}

From \texttt{off} with a saved pool (after \texttt{RestartFromDisk}), the CLI
starts playback at a chosen ready Part without generating anything. A full pool
enters \texttt{playingRotating}; a partial pool enters \texttt{playingFilling}
\emph{with the fill inactive} (\texttt{filling} false) --- this is the one place
the model relies on decoupling \texttt{filling} from the playing mode
(finding~\#2). A client may later re-arm the fill by claiming the Program with
\texttt{TurnOn}.

\begin{schema}{StartFromDisk}
  \Delta Program \\
  target? : PART
  \where
  mode = off \\
  target? \in pool \\
  format' = format \\
  pool' = pool \\
  failedParts' = failedParts \\
  playing' = \{ target? \} \\
  lastPlayed' = \emptyset \\
  attempts' = 0 \\
  lastError' = \emptyset \\
  filling' = zfalse \\
  (\# pool = (poolSize~format) \implies mode' = playingRotating) \\
  (\# pool < (poolSize~format) \implies mode' = playingFilling)
\end{schema}

\section{Consume-only Replay --- the Radio}
\label{sec:radio}

Alongside the single \texttt{Program} modelled above, Vox plays a second,
consume-only surface: a \emph{Radio} over a fixed \emph{Selection} of ready
Parts drawn from one or more Programs. Per-vibe pools (\texttt{vox-q7vh}) make
this the natural home for a \emph{union}: ``play trance calm'' may resolve to
two full albums at once. A Radio generates nothing, fills nothing, and ---
crucially --- carries \emph{no} \texttt{poolSize} bound, so a Selection may be
arbitrarily large. It is \texttt{Rotate} without the generation machinery.

A \emph{Selection} is a set of Parts. As with the file token behind a Part, the
per-track directory that resolves each Part is an implementation concern the
model keeps out of scope; a Selection is opaque set membership, nothing more.

\begin{zed}
  Selection == \power PART
\end{zed}

The \texttt{Radio} state carries only a Selection and the same ``optional Part''
cursor pair as a Program --- \texttt{playing} and \texttt{lastPlayed}, each a set
of at most one Part, both drawn from the Selection. It has \emph{no}
\texttt{pool}, \texttt{filling}, \texttt{attempts}, \texttt{failedParts},
\texttt{lastError}, or \texttt{mode}: a Radio neither generates nor fails, so it
needs none of the resilience state. Above all it carries no \texttt{poolSize}
bound.

The final conjunct, $selection \neq \emptyset \implies \# playing = 1$, ties the
cursor to the Selection: a non-empty Selection is an \emph{active} Radio and must
be playing exactly one Part, so the state schema forbids the ``active radio
playing nothing'' configuration ($selection \neq \emptyset$, $playing =
\emptyset$) outright rather than leaving it merely unreachable. With
$playing \subseteq selection$ and $\# playing \leq 1$ already in force, the
biconditional $\# playing = 1 \iff selection \neq \emptyset$ holds: the empty
Selection is the off state, where $playing = \emptyset$.

\begin{schema}{Radio}
  selection : Selection \\
  playing : \power PART \\
  lastPlayed : \power PART
  \where
  \# playing \leq 1 \land \# lastPlayed \leq 1 \\
  playing \subseteq selection \land lastPlayed \subseteq selection \\
  selection \neq \emptyset \implies \# playing = 1
\end{schema}

The missing bound is the whole point. A union of two full 12-Part albums is a
Selection of 24 Parts: a legal \texttt{Radio} state, though never a legal
\texttt{Program} state, whose invariant forces $\# pool \leq (poolSize~format)$.
This is precisely the delta \texttt{vox-q7vh} needs, and precisely why a union
cannot be shoehorned through \texttt{Program}.

\subsection{StartRadio}

Begin a Radio on a non-empty Selection, playing at one of its members. Nothing
has played yet, so \texttt{lastPlayed} is empty.

\begin{schema}{StartRadio}
  Radio' \\
  selection? : Selection
  \where
  selection? \neq \emptyset \\
  selection' = selection? \\
  playing' \subseteq selection' \land \# playing' = 1 \\
  lastPlayed' = \emptyset
\end{schema}

Retuning the writer from one Selection to another --- a
\texttt{SwitchSelection} --- is not a distinct
transition: it is \texttt{StartRadio} on the new Selection, exactly as
\texttt{SwitchProgram} is a fresh start on a new Program. Naming the alias keeps
the model vocabulary self-contained and realises the \texttt{SwitchProgram} /
\texttt{SwitchSelection} verb-parity.

\begin{zed}
  SwitchSelection \defs StartRadio
\end{zed}

\subsection{RadioRotate --- advance, next, and skip}

When the current Part ends, move to another member of the Selection, avoiding an
immediate repeat when the Selection holds more than one; a singleton Selection
replays. This is \texttt{Rotate} with every generation field removed --- there
is no \texttt{pool}, \texttt{filling}, \texttt{attempts}, \texttt{failedParts},
or \texttt{lastError} to frame, because a Radio has none.

\begin{schema}{RadioRotate}
  \Delta Radio
  \where
  \# selection \geq 1 \\
  \# playing = 1 \\
  selection' = selection \\
  lastPlayed' = playing \\
  (\# selection \geq 2 \implies \\
    \quad~ (playing' \subseteq selection \land \# playing' = 1 \\
    \quad~ \land playing' \neq playing)) \\
  (\# selection = 1 \implies playing' = playing)
\end{schema}

The precondition $\# playing = 1$ is redundant with the state invariant --- since
\texttt{RadioRotate} is $\Delta Radio$, the before-state is a valid \texttt{Radio},
and $selection \neq \emptyset$ (from $\# selection \geq 1$) already forces
$\# playing = 1$ --- but stating it makes plain that \texttt{RadioRotate}
\emph{advances} an already-playing Radio and never starts one: the empty-cursor
case is \texttt{StartRadio}, not a rotate.

As with \texttt{Rotate}, the automatic track-end advance and the user-driven
\texttt{next}/\texttt{skip} are one transition; a Radio has no owner and no
generation, so \texttt{RadioNext} is the only control it needs.

\begin{zed}
  RadioNext \defs RadioRotate
\end{zed}

\subsection{StaleFillIgnored --- the lost race}

A background fill belongs to the Program machine: it computes a fresh Part and
the single writer applies the outcome to whatever source is active. Suppose a
client has since retuned the writer to a Radio (a \texttt{SwitchSelection})
while a fill task was still in flight. The produced Part now arrives at a state
with no \texttt{filling} flag and no \texttt{pool} to grow. \texttt{FillOk}'s
success guard --- a fill in flight, $filling = ztrue$ --- has no inhabitant in
the \texttt{Radio} state, so the successful branch is disabled: its precondition
is false. Only the no-op branch remains, and it changes nothing.

\begin{schema}{StaleFillIgnored}
  \Xi Radio \\
  produced? : PART
\end{schema}

The schema is total in \texttt{produced?} and yields $\Xi Radio$: whatever the
fill delivered, a Radio consumes none of it. This is the model image of the
runtime lost race --- the writer narrows $isinstance(source, Program)$, the
narrowing fails against a Radio, and the writer logs the stale outcome at
\texttt{INFO} and survives. It is an \emph{enabled}, total $\Xi Radio$ no-op ---
the operation is defined for every \texttt{produced?} and simply changes nothing
--- never an error, and
never the \texttt{AttributeError} that a Radio without a \texttt{fill\_ok} would
otherwise raise (the \texttt{bas7}-class crash a transition-level model catches
at design time).

\subsection{RadioOff --- leaving the Radio}

A \texttt{TurnOff} while a Radio is active is a \emph{stop}, the exact parallel
of \texttt{TurnOff} on a Program: it tears the source down. The Selection is
emptied and the cursor cleared, leaving the empty Radio --- the model's
representation of ``no Radio active'', the same shape \texttt{StartRadio} refuses
as a starting Selection ($selection? \neq \emptyset$). It is a total operation on
the \texttt{Radio} schema, enabled from any Radio state, and is emphatically
\emph{not} a rejected lost race: leaving a Radio is user intent, always valid ---
the counterpart of the generate-family reject that \texttt{StaleFillIgnored}
models.

\begin{schema}{RadioOff}
  \Delta Radio
  \where
  selection' = \emptyset \\
  playing' = \emptyset \\
  lastPlayed' = \emptyset
\end{schema}

The empty Selection forces the cursor empty through the state invariant
($playing' \subseteq selection'$ with $selection' = \emptyset$), so clearing
\texttt{playing'} and \texttt{lastPlayed'} explicitly is a statement of intent
rather than a separate obligation. This mirrors \texttt{TurnOff} on a Program,
which clears \texttt{playing} to $\emptyset$ and drops \texttt{mode} to
\texttt{off}: a Radio has no \texttt{mode}, so its off state \emph{is} the empty
Selection.

\subsection{SwitchProgram --- source-agnostic retarget}

The channel that plays audio holds \emph{one} source at a time, and that source
is a union: either a Program or a Radio. A generate-mode start or switch ---
\texttt{StartProgram}/\texttt{SwitchProgram} in the design's vocabulary ---
retargets the channel to a fresh Program regardless of what it held before, so
it may fire while a Radio is active, transitioning Radio $\to$ Program. This is
the exact mirror of \texttt{StartRadio}/\texttt{SwitchSelection} retargeting
\emph{from} a Program to a Radio: a start or switch replaces whichever source is
current, never narrowing on its shape. Both directions are user intent, and both
are valid.

\texttt{SwitchProgram} discards the Radio and initialises a Program over the
resolved disk pool \texttt{diskPool?}, with the same pool-shape classification as
\texttt{TurnOn} --- an empty pool generates the first Part, a partial pool plays
and resumes the fill, a full pool plays and rotates. The Radio before-state is
unconstrained, so the retarget is enabled whether the Radio was playing or
already empty; that is what ``source-agnostic'' means, and it is why the before
half of $\Delta$ names \texttt{Radio} but the predicate never reads it.

\begin{schema}{SwitchProgram}
  Radio \\
  Program' \\
  fmt? : Format \\
  diskPool? : \power PART
  \where
  \# diskPool? \leq (poolSize~fmt?) \\
  format' = fmt? \\
  pool' = diskPool? \\
  failedParts' = \emptyset \\
  lastPlayed' = \emptyset \\
  attempts' = 0 \\
  lastError' = \emptyset \\
  (diskPool? = \emptyset \implies \\
    \quad~ (mode' = generatingFirst \land playing' = \emptyset \land filling' = ztrue)) \\
  (\# diskPool? \geq 1 \land \# diskPool? < (poolSize~fmt?) \implies \\
    \quad~ (mode' = playingFilling \land playing' \subseteq diskPool? \\
    \quad~ \land \# playing' = 1 \land filling' = ztrue)) \\
  (\# diskPool? = (poolSize~fmt?) \implies \\
    \quad~ (mode' = playingRotating \land playing' \subseteq diskPool? \\
    \quad~ \land \# playing' = 1 \land filling' = zfalse))
\end{schema}

The before-state \texttt{Radio} is present only to fix the source the transition
leaves; every primed field is determined by \texttt{diskPool?} alone, exactly as
\texttt{VibeStyleChange} determines a Program's next state from \texttt{newPool?}.
The symmetric retarget \emph{from} a Program to a Radio is not a Program schema
either: \texttt{StartRadio} produces a \texttt{Radio'} without naming the source
it displaced, because the displaced source is discarded whole. The model thus
keeps the two states apart and treats the channel-level union at the
orchestration seam, precisely as it abstracts the catalog lookup behind
\texttt{diskPool?} (finding~\#7).

\subsection{RadioVibeIgnored --- the vibe no-op}

A \texttt{VibeStyleChange} carries a new \emph{(vibe, style)} pool for the
generate machine. A Radio is consume-only: it has no vibe adaptation, and the
design locks that a drifting session vibe must never abandon a curated Selection.
A vibe change routed to an active Radio is therefore a deliberate no-op --- an
\emph{enabled}, total $\Xi Radio$ transition that changes nothing, of exactly
the shape of \texttt{StaleFillIgnored}, not an error.

\begin{schema}{RadioVibeIgnored}
  \Xi Radio \\
  newPool? : \power PART
\end{schema}

The schema is total in \texttt{newPool?} and yields $\Xi Radio$: whatever pool a
forwarded vibe offers, a Radio takes none of it. This is distinct from a lost
race --- nothing is ``lost'' --- and distinct from the generate-family reject: the
vibe simply does not apply to replay, so the ``no vibe-switch'' lock is a fact of
the model. With \texttt{RadioNext} $\defs$ \texttt{RadioRotate} already giving
\texttt{Next} on a Radio, the user-intent controls over a Radio are now complete:
\texttt{StartRadio}/\texttt{SwitchSelection} to enter, \texttt{RadioNext} to
advance, \texttt{SwitchProgram} to retarget to generate, and \texttt{RadioOff}
to stop.

\section{Key Properties}

The following hold by construction (state invariants) or as operation
postconditions; they are the intended proof obligations for animation.

\begin{itemize}
  \item \textbf{Bounded pool}: $\# pool \leq (poolSize~format)$ always; for
    \texttt{playlist}, $\# pool \leq 12$.
  \item \textbf{Generation only below full}: \texttt{FillOk} requires
    $\# pool < (poolSize~format)$; at $\# pool = (poolSize~format)$ the mode is
    \texttt{playingRotating} with $filling = zfalse$, so no fill runs.
  \item \textbf{At most one fill}: \texttt{filling} is a single flag, and
    $filling = ztrue$ implies \texttt{mode} is \texttt{generatingFirst} or
    \texttt{playingFilling}.
  \item \textbf{Playing is ready}: whenever $\# playing = 1$,
    $playing \subseteq pool$, and every member of \texttt{pool} is a
    \texttt{ready} Part; so the playing Part is always ready.
  \item \textbf{Failed is observable}: $mode = failed \implies \# lastError = 1$,
    and the two operations that reach \texttt{failed}
    (\texttt{FirstTrackBadPrompt}, \texttt{RetryExhausted}) set both
    \texttt{mode} and \texttt{lastError} --- there is no silent disable.
  \item \textbf{Full pool never hard-fails}: \texttt{failed} requires
    $pool = \emptyset$, and no operation empties a non-empty pool, so a
    Program with $\# pool \geq 1$ can never reach \texttt{failed}. In
    particular a full ($\# pool = poolSize$) pool rotates from disk forever.
  \item \textbf{Two failure surfaces}: a \emph{program-level} failure sets
    \texttt{lastError} (in \texttt{retrying}/\texttt{failed}); a
    \emph{per-Part} permanent failure is recorded in \texttt{failedParts} while
    the Program stays healthy (\texttt{FillBadPart}). Both are queryable; the
    log is never a client interface.
  \item \textbf{Replay generates nothing}: \texttt{PlayPart}, \texttt{Loop},
    \texttt{Rotate}, \texttt{StartFromDisk} all leave \texttt{pool} and
    \texttt{failedParts} unchanged and never set \texttt{filling} true; their
    precondition forces the target to be a \texttt{ready} Part. The CLI can
    drive them with no LLM and zero credits.
  \item \textbf{No immediate repeat}: \texttt{Rotate} with $\# pool \geq 2$
    gives $playing' \neq playing$; \texttt{PlayPart} deliberately does not, as
    it honours an explicit index.
  \item \textbf{Retry cap bounds only the empty-pool case}:
    \texttt{RetryExhausted} requires $pool = \emptyset$, so the cap forces
    \texttt{failed} only when nothing is playing; a non-empty pool tolerates
    transient errors indefinitely, recovering via \texttt{Recover} while
    playback continues.
  \item \textbf{Off clears generation state}: $mode = off \implies$
    \texttt{playing}, \texttt{filling}, \texttt{lastError}, and
    \texttt{failedParts} are all cleared, so a stopped Program carries no stale
    fill, error, or per-Part failure into its next \texttt{TurnOn}.
  \item \textbf{No session gate}: no operation takes or checks a session. Every
    transition is enabled on \texttt{mode} (and pool shape) alone, so the same
    command --- \texttt{next}, \texttt{skip}, \texttt{vibe}, \texttt{off} ---
    succeeds from any client or process, and the recovery path
    \texttt{VibeStyleChange} is always available from \texttt{failed}.
  \item \textbf{Replay generates nothing and is uncapped}: \texttt{StartRadio}
    and \texttt{RadioRotate} touch no generation field --- a \texttt{Radio} has
    none --- and the \texttt{Radio} state carries no \texttt{poolSize} bound, so
    a Selection may be arbitrarily large. A union of two full albums (24 Parts)
    is a legal \texttt{Radio} state though never a legal \texttt{Program} state.
  \item \textbf{A stale fill against a Radio never mutates}: \texttt{FillOk}'s
    success guard ($filling = ztrue$) is false on a \texttt{Radio}, so a Part
    produced by an in-flight fill and routed to an active Radio leaves it
    unchanged (\texttt{StaleFillIgnored}, $\Xi Radio$) --- never a crash, never
    a corrupt cursor.
  \item \textbf{Leaving a Radio is a valid stop}: \texttt{RadioOff} empties the
    Selection and cursor from any Radio state ($\Delta Radio$ with
    $selection' = \emptyset$, whence $playing' = lastPlayed' = \emptyset$ by the
    invariant), the exact parallel of \texttt{TurnOff} on a Program. It is user
    intent, always enabled, never a rejected lost race.
  \item \textbf{Retarget is source-agnostic}: the user-intent signals --- off,
    on, switch, next --- act on \emph{either} source. \texttt{SwitchProgram}
    retargets an active Radio to a Program ($Radio \to Program'$), the mirror of
    \texttt{StartRadio}/\texttt{SwitchSelection} retargeting a Program to a
    Radio; \texttt{RadioNext} advances; \texttt{RadioOff} stops. Only the
    generate-internal outcome (\texttt{StaleFillIgnored}) and the auto
    vibe-switch (\texttt{RadioVibeIgnored}) are disabled against a Radio, and
    both as $\Xi Radio$ no-ops rather than errors.
\end{itemize}

\section{Findings from Formalisation}
\label{sec:findings}

Writing the model surfaced these decisions, which the Phase-1 code design must
realise consistently. Each is a place where the informal concept was silent,
underspecified, or self-contradictory.

\begin{enumerate}
  \item \textbf{The six modes are format-general; only \texttt{poolSize} and
    the end-of-list behaviour differ.} \texttt{playingRotating} means ``full
    pool'' for every format --- looping for playlist, ``all Parts generated,
    playing to the end'' for podcast/audiobook. The loop-versus-stop choice
    lives in \texttt{Rotate} (which for a finite format would additionally test
    ``is this the last Part''), not in a new mode. This is what lets podcast and
    audiobook extend the model by supplying a \texttt{poolSize} and a
    format-specific \texttt{Rotate}, not a new state machine.

  \item \textbf{\texttt{filling} must be decoupled from \texttt{playingFilling}.}
    The old four-state model conflated ``pool not yet full'' with ``a fill is
    running''. The moment the CLI can play a partial pool without generating
    (\texttt{StartFromDisk}), and the moment a transient error pauses the fill
    while a partial pool keeps playing, that conflation is contradictory: a
    playing partial pool with no active fill has no mode. The model therefore
    makes \texttt{playingFilling} mean ``playing, pool $<$ full'' and lets
    \texttt{filling} be an orthogonal flag constrained by
    $filling = ztrue \implies mode \in \{generatingFirst, playingFilling\}$.
    Without this decoupling the CLI cold-start and the transient-recovery paths
    cannot be expressed.

  \item \textbf{Playback must remain enabled in \texttt{retrying}.} A transient
    backoff pauses \emph{generation}, not \emph{playback}. \texttt{Rotate},
    \texttt{Loop}, and \texttt{PlayPart} are therefore enabled in
    \texttt{retrying} (with a non-empty pool); only \texttt{FillOk} is not.
    The original model, which lacked \texttt{retrying}, enabled advance only in
    the two playing modes; the resilience generalisation had to widen that
    guard, or a track ending during backoff would have no legal transition.

  \item \textbf{Hard-failure is confined to the empty-pool case, structurally.}
    \texttt{failed} carries $pool = \emptyset$ in the state invariant, and the
    only two operations that reach it require an empty pool
    (\texttt{FirstTrackBadPrompt} from \texttt{generatingFirst};
    \texttt{RetryExhausted} which tests $pool = \emptyset$). A transient error
    on a non-empty pool (\texttt{FillTransient}) enters \texttt{retrying} but can
    only leave via \texttt{Recover} --- never \texttt{RetryExhausted} --- so
    ``a full pool never hard-fails'' (\texttt{vox-ig52} acceptance criterion~3)
    is a theorem of the state schema, not a runtime check. The corner the
    concept left open --- ``what happens when fill on a non-empty pool exhausts
    its retries?'' --- is answered by \texttt{RetryCapped}: it does not exhaust;
    it plays on and keeps trying at the capped backoff, self-looping in
    \texttt{retrying} with \texttt{attempts} pinned at the cap.

  \item \textbf{There are two failure surfaces, and both must be observable.}
    A \emph{program-level} failure (nothing can play) sets \texttt{lastError}
    and shows as \texttt{mode = failed} or \texttt{retrying}. A \emph{per-Part}
    permanent failure while the Program plays on (\texttt{FillBadPart}) is
    recorded in \texttt{failedParts} with a reason, and the program-level
    \texttt{lastError} stays empty because the Program is healthy. The
    invariant $\# lastError = 1 \implies mode \in \{retrying, failed\}$ forbids a
    healthy playing Program from carrying a program-level error, which forced
    the per-Part surface to exist. \texttt{vox-ig52}'s silent-failure defect is
    closed only if \emph{both} surfaces are exposed through the \texttt{status}
    tool; a design that surfaces only \texttt{lastError} would silently drop the
    \texttt{FillBadPart} case.

  \item \textbf{No operation requires ownership --- there is no ownership.}
    \emph{(Superseded by the operator ruling of 2026-07-05; retained as errata.)}
    An earlier draft distinguished agent operations, which owned a Program, from
    consume-only operations, which did not, and had to accept the asymmetric
    consequence that an agent could not \texttt{Skip} or \texttt{VibeStyleChange}
    a CLI-started, unowned Program without first claiming it. The operator ruled
    ownership out entirely: it is complex, bug-ridden, and unusable across
    sessions --- gating control on the turning-on session made music
    uncontrollable from any other session. The \texttt{SESSION} type, the
    \texttt{owner} state variable, the \texttt{who?} inputs, and every ownership
    invariant are gone. Every operation is now enabled on \texttt{mode} (and
    pool shape) alone; \texttt{voxd} state is machine-universal, and any client
    from any process drives every command. The old worry this finding recorded
    --- ``can the CLI play a pool it did not turn on?'' --- is answered
    unconditionally yes.

  \item \textbf{Part indexing is a resolved surface parameter, not modelled
    state.} \texttt{PlayPart} takes a Part, with precondition
    $target? \in pool$; the CLI's 1-based index (\texttt{playlist:2}) is
    resolved to a Part against a stable numbering of \texttt{pool} before the
    operation runs. The model deliberately abstracts the numbering away: the
    property worth proving is ``the addressed Part must be ready and nothing is
    generated'', and that is exactly the precondition plus the framing. The
    design owns the ordering (concept: Programs are an \emph{ordered} list); an
    out-of-range index is a resolution failure the CLI reports before any
    operation is attempted, so no ``index out of range'' transition appears
    here. The same abstraction covers \emph{which} pool a \texttt{TurnOn} or
    \texttt{VibeStyleChange} enters: with per-vibe pools (\texttt{vox-q7vh}) the
    \texttt{diskPool?}/\texttt{newPool?} it binds is resolved by a catalog tag
    or name lookup before the operation runs, exactly as the 1-based index is
    resolved to a Part, so the model abstracts the lookup away and no
    \texttt{Program} schema changes to accommodate per-vibe pools.

  \item \textbf{\texttt{failed} is recoverable, and recovery is unconditional.}
    \emph{(The ownership half of this finding is moot after the 2026-07-05
    ruling; the recoverability half stands.)} \texttt{failed} is not terminal:
    it is a queryable state that \texttt{VibeStyleChange} re-authors out of, so
    no design may treat it as a permanent disable. The earlier draft added a
    complication here --- \texttt{failed} had to \emph{retain} its owner, or no
    session could retune it and the failure would be unrecoverable except by a
    full \texttt{TurnOff}/\texttt{TurnOn}. With ownership removed that concern
    evaporates: \texttt{VibeStyleChange} is gated only on $mode \neq off$, so any
    client can retune a \texttt{failed} Program back into generation. The
    recovery path is strictly \emph{more} available than the owned model made
    it.

  \item \textbf{The atomic-delivery abstraction hides \texttt{pending} and
    \texttt{generating}.} Phase-1 generation is modelled as an atomic step
    (\texttt{FillOk} produces a ready Part; there is no observable half-generated
    Part), so per-Part \texttt{pending}/\texttt{generating} statuses collapse
    into the single \texttt{filling} flag. This is sound for music, whose Parts
    are short, but the spoken formats generate for minutes; Phase~2/3 will need
    to promote those statuses to stored per-Part state and add the corresponding
    transitions. The \texttt{PartStatus} free type is declared now so that
    extension does not disturb the state signature.
\end{enumerate}

\end{document}
